\documentclass[11pt]{article}

\usepackage[utf8]{inputenc}
\usepackage[T1]{fontenc}
\usepackage{geometry}
\usepackage{amsmath}
\usepackage{amssymb}
\usepackage{booktabs}
\usepackage{hyperref}
\usepackage{xurl}
\usepackage{xcolor}
\usepackage{listings}
\usepackage{enumitem}
\usepackage{parskip}
\usepackage{graphicx}
\graphicspath{{../images/}}

\geometry{margin=1in}

\lstset{
  language=Java,
  basicstyle=\ttfamily\small,
  keywordstyle=\color{blue},
  commentstyle=\color{gray},
  stringstyle=\color{teal},
  showstringspaces=false,
  columns=fullflexible,
  frame=single,
  framerule=0.2pt,
  breaklines=true
}

\setlist[itemize]{topsep=0.3em,itemsep=0.2em}
\setlist[enumerate]{topsep=0.3em,itemsep=0.2em}

% Simple per-section source line.
\newcommand{\notesources}[1]{\par\noindent{\small\color{gray}\textit{Sources:} \sloppy #1\par}}

% Unit-wise source strings for this subject (used by slides/notes).
% Path is relative to the lecture `latex/` folder (the compilation CWD).
% OOPS with Java (CSEG1044) — unit-wise sources
%
% These are short, student-facing reference strings intended to be used with:
%   - \slidesources{...}  (in beamer slides)
%   - \notesources{...}   (in lecture notes)
%
% Style inspiration: other subjects in My-Subjects/ use semicolon-separated sources
% and include an "accessed" date for web documentation.

\newcommand{\OOPSUnitISources}{Schildt, \textit{Java: A Beginner's Guide} (10e), McGraw-Hill (Java fundamentals + OOP basics).; Oracle Java Tutorials: Getting Started + Language Basics + Classes and Objects (accessed \today).; Java SE API docs (e.g., \texttt{String}, \texttt{Scanner}) (accessed \today).}

\newcommand{\OOPSUnitIISources}{Schildt, \textit{Java: The Complete Reference} (12e), McGraw-Hill (classes, inheritance, packages, interfaces).; Bloch, \textit{Effective Java} (3e), Addison-Wesley, 2018 (design choices; interfaces vs abstract classes).; Oracle Java Tutorials: Classes and Objects + Interfaces + Packages (accessed \today).; Java SE API docs (e.g., \texttt{Comparable}, \texttt{Runnable}) (accessed \today).}

\newcommand{\OOPSUnitIIISources}{Schildt, \textit{Java: The Complete Reference} (12e) (nested/inner classes, exceptions, threads, I/O).; Oracle Java Tutorials: Nested Classes + Exceptions + Concurrency + Basic I/O (accessed \today).; Java SE API docs (e.g., \texttt{Thread}, \texttt{AutoCloseable}) (accessed \today).}

\newcommand{\OOPSUnitIVSources}{Schildt, \textit{Java: The Complete Reference} (12e) (generics, lambdas, Swing, JDBC).; Bloch, \textit{Effective Java} (3e) (generics + lambdas).; Oracle Java Tutorials: Generics + Lambda Expressions + Swing + JDBC Basics (accessed \today).; Java SE API docs (e.g., \texttt{java.util.function}, \texttt{javax.swing}, \texttt{java.sql}) (accessed \today).}

\newcommand{\OOPSUnitVSources}{Schildt, \textit{Java: The Complete Reference} (12e) (Collections Framework + wrapper classes).; Oracle Java docs: Collections Framework overview (accessed \today).; Java SE API docs (\texttt{java.util} collections; wrapper classes in \texttt{java.lang}) (accessed \today).}

\newcommand{\OOPSUnitVISources}{Git SCM: \textit{Pro Git} (2e) (version control workflow), accessed \today.; GitHub Docs (repositories, issues, pull requests), accessed \today.; Oracle Java Tutorials: Swing + JDBC (accessed \today).; JUnit 5 User Guide (accessed \today).; Apache Maven Project: Getting Started (accessed \today).; Gradle User Manual: Getting Started (accessed \today).; UML class diagrams: UML 2.x overview/spec (accessed \today).}

\newcommand{\OOPSMiscWhyInterfacesSources}{Schildt, \textit{Java: The Complete Reference} (12e) (interfaces).; Bloch, \textit{Effective Java} (3e) (prefer interfaces where appropriate).; Oracle Java Tutorials: Interfaces (accessed \today).; Java SE API docs (e.g., \texttt{Comparable}, \texttt{Runnable}) (accessed \today).}



\title{Object Oriented Programming (CSEG1044)\\Unit 06 -- Lecture 02 Notes\\Requirements → Use Cases → Milestones}
\author{Tofik Ali}
\date{\today}

\begin{document}
\maketitle
\tableofcontents

\section{Lecture Overview}
In Deliverable-0 you selected a theme and defined a scope.
Now we convert that scope into \textbf{clear requirements} that can be implemented and tested.
The workflow is:
\begin{center}
Requirements $\rightarrow$ Use Cases $\rightarrow$ Milestones/Backlog
\end{center}
\notesources{\OOPSUnitVISources}

\section{Syllabus Alignment (Unit 06)}
\textbf{Unit 06:} requirements + use cases + project planning

\section{Learning Outcomes}
After this lecture, you should be able to:
\begin{itemize}
  \item Identify actors and write use cases with success/failure scenarios and acceptance criteria.
  \item Create Deliverable-1: use case diagram + 3 written use cases.
  \item Convert requirements into milestones and update your GitHub backlog.
\end{itemize}

\section{Key Terms (Quick)}
\begin{itemize}
  \item \textbf{Requirement}: what the system must do (functional) or must be (non-functional).
  \item \textbf{Functional requirement}: a feature/behavior (e.g., create a booking).
  \item \textbf{Non-functional requirement}: quality/constraint (e.g., usability, security, performance).
  \item \textbf{Actor}: external role interacting with the system (admin, user).
  \item \textbf{Use case}: goal-oriented interaction between an actor and the system.
  \item \textbf{Scenario}: step-by-step story of using the system (success/failure paths).
  \item \textbf{Acceptance criteria}: conditions to consider a use case complete (testable).
  \item \textbf{Definition of done (DoD)}: your team's shared checklist for when work is truly ``complete''.
  \item \textbf{Milestone}: group of backlog items delivered by a date/session.
\end{itemize}

\section{Detailed Notes}
\subsection{Step 1: Convert scope into requirements}
Your Deliverable-0 scope (Must/Should/Could/Non-goals) is still high level.
Now we convert it into requirements that can be \textbf{implemented} and \textbf{tested}.

\subsubsection{How to write a good requirement}
Good requirements are:
\begin{itemize}
  \item \textbf{clear:} avoid vague words like ``fast'' or ``good'',
  \item \textbf{specific:} mention the feature and the rule,
  \item \textbf{testable:} you can write 2--3 test cases from it.
\end{itemize}

\textbf{Weak (avoid):}
\begin{itemize}
  \item ``The system should be fast.''
  \item ``The UI should be user friendly.''
\end{itemize}

\textbf{Better (testable):}
\begin{itemize}
  \item ``Booking cannot overlap an existing booking for the same resource.''
  \item ``Capacity must be a positive integer; otherwise show an error message and do not save.''
  \item ``On success, the record is saved in DB and appears in the list view without restarting the app.''
\end{itemize}

\subsection{Step 2: Identify actors and use cases (goals)}
Start from your requirements and identify:
\begin{itemize}
  \item actors (roles): Admin, User, Staff, etc.,
  \item goals: what each actor wants to achieve.
\end{itemize}
Each goal becomes a use case.

\textbf{Example (Campus Resource Booking):}
\begin{itemize}
  \item Actor: Admin $\rightarrow$ Add/Update Resource
  \item Actor: User $\rightarrow$ Create Booking, Cancel Booking, View Availability
\end{itemize}

\subsection{Step 3: Create a use case diagram (quick guide)}
The use case diagram is a \textbf{map} of your features:
\begin{itemize}
  \item draw a system boundary rectangle (your application),
  \item put actors outside (roles),
  \item put use cases inside as verb phrases (``Create Booking'', ``Add Resource'').
\end{itemize}
\textbf{Tip:} keep the diagram simple; use \texttt{include/extend} only if you truly need it.

\subsection{Step 4: Write the use cases (template)}
Write each use case in a structured way. Keep it short but complete:
\begin{itemize}
  \item Use case name
  \item Primary actor
  \item Preconditions
  \item Trigger
  \item Main success scenario (numbered steps)
  \item Alternate flows (failures/edge cases)
  \item Postconditions
  \item Acceptance criteria (3--6 testable bullets)
\end{itemize}

\subsection{Worked example: Use case ``Create Booking'' (full)}
\textbf{Use case name:} Create Booking\par
\textbf{Primary actor:} User\par
\textbf{Preconditions (examples):}
\begin{itemize}
  \item At least one resource exists in the system.
  \item Database connection is available.
\end{itemize}
\textbf{Trigger:} user clicks ``Create Booking'' / ``Save Booking'' button.

\subsubsection*{Main success scenario (example)}
\begin{enumerate}
  \item System shows booking form (resource, start time, end time).
  \item User selects a resource and enters start/end time.
  \item System validates input (mandatory fields, start $<$ end).
  \item System checks overlap for the same resource.
  \item System saves booking in DB.
  \item System shows success message.
  \item System updates booking list view (the new booking is visible).
\end{enumerate}

\subsubsection*{Alternate flows (examples)}
\begin{itemize}
  \item \textbf{A1: Invalid time range} (start $\ge$ end)\par
  System shows validation error and does not save.

  \item \textbf{A2: Overlap exists}\par
  System shows ``Overlapping booking'' message and does not save.

  \item \textbf{A3: DB failure}\par
  System shows a friendly error message and does not crash.
\end{itemize}

\subsubsection*{Postconditions (examples)}
\begin{itemize}
  \item On success, booking record exists in DB and is visible in UI list.
  \item On failure, no record is created (data remains consistent).
\end{itemize}

\subsubsection*{Acceptance criteria (sample)}
\begin{itemize}
  \item Booking cannot overlap an existing booking for the same resource.
  \item Boundary case: end time equal to another booking's start time is allowed.
  \item Invalid input shows a clear message near the form and does not save.
  \item On success, booking is saved and visible in list view.
\end{itemize}

\subsection{Step 5: Convert use cases into backlog items}
One use case usually becomes \textbf{multiple issues} (UI + service + DAO + DB).
\textbf{Example breakdown for ``Create Booking'':}
\begin{itemize}
  \item UI: create booking form panel + input validation
  \item Service: \texttt{BookingService.createBooking(...)} with overlap rules
  \item DAO: \texttt{BookingDao.insert(...)} + query to find overlaps
  \item DB: schema constraints + indexes (by \texttt{resource\_id}, time range)
  \item Test/demo: 3 scenarios (success, overlap, invalid time)
\end{itemize}

\subsection{Step 6: Milestones (turn backlog into a plan)}
Milestones help you avoid ``everything at the end''.
Keep milestones aligned with Unit 06 deliverables:
\begin{itemize}
  \item \textbf{Milestone-1 (Deliverable-1):} diagram + 3 written use cases
  \item \textbf{Milestone-2 (Deliverable-2):} class diagram + package structure
  \item \textbf{Milestone-3 (Deliverable-3):} DB schema + DAO CRUD for 1--2 entities
  \item \textbf{Milestone-4 (Deliverable-4):} UI integration + validations + end-to-end flow
  \item \textbf{Milestone-5 (final):} testing, packaging, polish, demo readiness
\end{itemize}
\notesources{\OOPSUnitVISources}

\section{Deliverable-1: Use case diagram + 3 written use cases}
\subsection*{Deliverable-1A: Use case diagram}
Diagram should include:
\begin{itemize}
  \item 2--3 actors,
  \item 6--10 use cases,
  \item relationships if needed (include/extend --- optional).
\end{itemize}

\subsection*{Deliverable-1B: 3 written use cases (suggested)}
Pick the most important core flows:
\begin{itemize}
  \item Create Booking (core)
  \item Add Resource (admin)
  \item List/Search Bookings (read flow)
\end{itemize}

\subsection*{Minimum quality bar (what we will check)}
\begin{itemize}
  \item Each written use case has at least 6--10 steps in main success scenario.
  \item Each use case includes at least 2 alternate flows (invalid input + business rule failure).
  \item Acceptance criteria are testable and can be demonstrated during the final demo.
\end{itemize}

\section{In-Class Activity (evidence)}
\begin{itemize}
  \item Add Deliverable-1 document to repo (PDF/Doc) or link in README.
  \item Create/update GitHub issues based on use cases (one issue per use case / feature).
  \item Add milestone labels or GitHub milestones (optional).
\end{itemize}

\section{In-Class Exercises (with brief solutions)}
\begin{enumerate}
  \item Convert one Must-have feature into 2 requirements.\par
  \textbf{Sample:} ``Create booking'' $\rightarrow$ (1) booking must store resource + start + end, (2) booking must not overlap existing booking for same resource.

  \item Identify 2--3 actors for your project.\par
  \textbf{Sample:} Admin, User/Student/Faculty.

  \item Write 3 acceptance criteria for ``Add Resource''.\par
  \textbf{Sample:} name required; capacity positive; on success resource appears in list.
\end{enumerate}

\section{Self-Study Practice (no solutions)}
\begin{enumerate}
  \item Write a use case diagram for your project (6--10 use cases).
  \item Write 2 complete written use cases using the template (include alternate flows).
  \item Convert one use case into at least 5 GitHub issues (UI/service/DAO/DB/test).
\end{enumerate}

\section{Checkpoint Questions (with sample answers)}
\textbf{Q1.} Why do we write acceptance criteria?\par
\textbf{Sample answer:} Acceptance criteria make requirements testable. They define what ``done'' means, reduce confusion, and help validate the system during demo.

\vspace{0.6em}
\textbf{Q2.} What is the difference between an actor and a use case?\par
\textbf{Sample answer:} An actor is a role interacting with the system. A use case is a goal/action the actor performs using the system.

\vspace{0.6em}
\textbf{Q3.} Why do we separate ``main success scenario'' and ``alternate flows''?\par
\textbf{Sample answer:} It forces you to think about errors and edge cases early. Alternate flows often become validations and exception-handling code.

\end{document}
